\chapter{预测方法与示例}
\section{预测模型}
\subsection{滑动平均}
设 $x_t$ 为时刻 $t$ 的观测位移。长度为 $k$ 的滑动平均模型用最近 $k$ 个观测估计下一时刻的值，见式~\eqref{eq:average}。
\begin{equation}
  \widehat{x}_{t+1}=\frac{1}{k}\sum_{j=0}^{k-1}x_{t-j}.
  \label{eq:average}
\end{equation}

\section{图表}
表~\ref{tab:samples} 给出三条示例观测。图~\ref{fig:flow} 给出从观测序列到预测值的计算流程。

\begin{table}[htbp]
  \caption{位移观测示例}
  \label{tab:samples}
  \begin{tabularx}{\textwidth}{@{}*{2}{>{\centering\arraybackslash}X}@{}}
    \toprule
    采样时刻 $t$（\unit{\second}） & 观测位移 $x$（\unit{\milli\metre}）\\
    \midrule
    0 & 1.12\\
    1 & 1.25\\
    2 & 1.31\\
    \bottomrule
  \end{tabularx}
\end{table}

\begin{figure}[htbp]
  % 可用 \includegraphics[width=.75\textwidth]{your-figure.pdf} 替换示例图。
  \fbox{\strut 观测序列}\quad$\longrightarrow$\quad
  \fbox{\strut 滑动平均}\quad$\longrightarrow$\quad
  \fbox{\strut 预测值}
  \caption{预测计算流程}
  \label{fig:flow}
\end{figure}
